% ROZDZIAŁ 6 — plik włączany przez main.tex
\chapter{Funkcje mierzalne i ich aproksymacja}
\label{ch:funkcje-mierzalne}

\section{Funkcje jako podstawowy obiekt teorii miary}

Dotychczas mierzyliśmy zbiory. W praktyce chcemy mierzyć
również zbiory postaci $\{x:f(x)>a\}$, przybliżać
funkcje funkcjami o skończenie wielu wartościach, a następnie
określić całkę. Właśnie dlatego potrzebujemy pojęcia
funkcji mierzalnej. Ważne, by odróżniać \emph{mierzalność}
od \emph{całkowalności}: pierwsza pozwala dopiero
postawić drugie pytanie.

W całym rozdziale $(X,\mathcal F)$ oznacza przestrzeń mierzalną.
Jeśli potrzebna jest miara, zapisujemy $(X,\mathcal F,\mu)$.
Przez $\overline{\mathbb R}=[-\infty,\infty]$ rozumiemy
rozszerzoną prostą rzeczywistą wyposażoną w topologię
porządkową; jej $\sigma$-algebrą jest
$\mathcal B(\overline{\mathbb R})$. Każda funkcja
$\mathbb R$-wartościowa może być traktowana również jako
$\overline{\mathbb R}$-wartościowa.

\section{Funkcje rzeczywiste i rozszerzone rzeczywiste}

\begin{definicja}[Mierzalność funkcji rozszerzonej]
\label{def:extended-measurable}
Funkcję $f:X\to\overline{\mathbb R}$ nazywamy
\emph{mierzalną}, jeśli jest
$(\mathcal F,\mathcal B(\overline{\mathbb R}))$-mierzalna,
czyli $f^{-1}(B)\in\mathcal F$ dla każdego
$B\in\mathcal B(\overline{\mathbb R})$.
\end{definicja}

\begin{twierdzenie}[Testy mierzalności na półprostych]
\label{thm:extended-tests}
Niech $(X,\mathcal F)$ będzie przestrzenią mierzalną,
a $f:X\to\overline{\mathbb R}$. Następujące warunki są równoważne:
\begin{enumerate}[label=\textup{(\roman*)}]
\item $f$ jest mierzalna;
\item $\{f<a\}\in\mathcal F$ dla każdego $a\in\mathbb R$;
\item $\{f\le a\}\in\mathcal F$ dla każdego $a\in\mathbb R$;
\item $\{f>a\}\in\mathcal F$ dla każdego $a\in\mathbb R$;
\item $\{f\ge a\}\in\mathcal F$ dla każdego $a\in\mathbb R$.
\end{enumerate}
Wystarczy w każdym z warunków (ii)--(v) badać liczby wymierne.
\end{twierdzenie}
\begin{proof}
Zbiory postaci $[-\infty,a)$, $[-\infty,a]$,
$(a,\infty]$ oraz $[a,\infty]$ dla $a\in\mathbb Q$
generują tę samą $\sigma$-algebrę
$\mathcal B(\overline{\mathbb R})$. Istotnie, każda
z tych rodzin jest rozdzielona od pozostałych za pomocą
dopełnień i przeliczalnych granic. Na przykład
$[-\infty,a]=\bigcap_{n\ge1}[-\infty,a+1/n)$ oraz
$(a,\infty]=\overline{\mathbb R}\setminus[-\infty,a]$.
Każdy otwarty przedział w rozszerzonej prostej
jest przecięciem dwóch odpowiednich półprostych,
a każdy zbiór otwarty jest przeliczalną sumą takich
przedziałów, ponieważ $\overline{\mathbb R}$ ma bazę
przeliczalną. Teraz stosujemy twierdzenie
\ref{thm:measurable-generator}.
\end{proof}

\begin{uwaga}[Wartości nieskończone]
Dla mierzalnej funkcji $f:X\to\overline{\mathbb R}$
zbiory $\{f=+\infty\}=\bigcap_{n\ge1}\{f>n\}$ i
$\{f=-\infty\}=\bigcap_{n\ge1}\{f<-n\}$ są mierzalne.
Przy działaniach arytmetycznych musimy unikać
niezdefiniowanych wyrażeń $+\infty-\infty$,
$\infty/\infty$ lub $0\cdot\infty$, chyba że wcześniej
jawnie określimy przyjętą konwencję.
\end{uwaga}

\begin{przyklad}[Funkcja charakterystyczna zbioru]
\label{ex:indicator-measurable}
Dla zbioru $A\subseteq X$ definiujemy
$\chi_A(x)=1$ dla $x\in A$ oraz
$\chi_A(x)=0$ dla $x\notin A$.
Funkcja $\chi_A:X\to\mathbb R$ jest mierzalna
wtedy i tylko wtedy, gdy $A\in\mathcal F$.
Kierunek konieczny otrzymujemy z
$A=\chi_A^{-1}((1/2,\infty))$, a dostateczność
z faktu, że każdy przeciwobraz zbioru borelowskiego
przez $\chi_A$ jest jednym z czterech zbiorów:
$\varnothing$, $A$, $X\setminus A$, $X$.
\end{przyklad}

\section{Działania na funkcjach mierzalnych}

\begin{twierdzenie}[Działania algebraiczne]
\label{thm:measurable-algebra}
Niech $(X,\mathcal F)$ będzie przestrzenią mierzalną,
a $f,g:X\to\mathbb R$ funkcjami mierzalnymi.
Wtedy funkcje $f+g$, $f-g$, $fg$, $|f|$,
$\max(f,g)$ i $\min(f,g)$ są mierzalne.
Zbiór $\{g\ne0\}$ należy do $\mathcal F$, a
funkcja $f/g:\{g\ne0\}\to\mathbb R$ jest mierzalna
względem śladu $\sigma$-algebry. Jeżeli $g(x)\ne0$
dla każdego $x\in X$, to $f/g$ jest mierzalna na całym $X$.
\end{twierdzenie}
\begin{proof}
Odwzorowanie $(f,g):X\to\mathbb R^2$ jest mierzalne
na mocy wniosku \ref{cor:coordinate-measurable}.
Dodawanie $S(u,v)=u+v$, odejmowanie $D(u,v)=u-v$,
mnożenie $M(u,v)=uv$, maksimum i minimum są ciągłymi
funkcjami na $\mathbb R^2$. Moduł jest funkcją ciągłą
na $\mathbb R$. Ich złożenia z funkcją $(f,g)$ są
więc mierzalne na mocy twierdzenia
\ref{thm:composition-measurable}.
Zbiór $\{g\ne0\}=g^{-1}(\mathbb R\setminus\{0\})$
jest mierzalny; ograniczenie $(f,g)$ do niego jest
mierzalne. Funkcja $(u,v)\mapsto u/v$ jest ciągła
na otwartym zbiorze $\{(u,v):v\ne0\}$, co kończy dowód.
\end{proof}

\begin{twierdzenie}[Operacje na funkcjach rozszerzonych]
\label{thm:extended-operations}
Niech $(X,\mathcal F)$ będzie przestrzenią mierzalną,
a $f,g:X\to\overline{\mathbb R}$ funkcjami mierzalnymi.
Wtedy $-f$, $|f|$, $\max(f,g)$ oraz $\min(f,g)$
są mierzalne. Jeśli $f,g$ są nieujemne, ich suma
$f+g$ z naturalną regułą $a+\infty=\infty$
jest mierzalna. Jeżeli nie zachodzi w żadnym punkcie
jednocześnie $f=+\infty$, $g=-\infty$ ani
$f=-\infty$, $g=+\infty$, to również $f+g$ jest mierzalna.
\end{twierdzenie}
\begin{proof}
Negacja i moduł są ciągłe w topologii rozszerzonej
prostej. Funkcje maksimum i minimum są ciągłe na
$\overline{\mathbb R}^2$ z topologią produktową:
na przykład $\max(u,v)<a$ jest równoważne
$u<a$ i $v<a$, zaś $\max(u,v)>a$ zachodzi,
gdy $u>a$ lub $v>a$. Te testy wystarczają do
mierzalności z twierdzenia \ref{thm:extended-tests}.
Dla nieujemnych $f,g$ i dowolnego $a\in\mathbb R$
zachodzi
\[
\{f+g>a\}=\bigcup_{q\in\mathbb Q}
\bigl(\{f>q\}\cap\{g>a-q\}\bigr).
\]
Jeżeli $f(x)+g(x)>a$, można wybrać wymierne $q$
ściśle pomiędzy $a-g(x)$ i $f(x)$; gdy którakolwiek
wartość wynosi $+\infty$, odpowiedni przedział
wymiernych nadal jest niepusty. Implikacja odwrotna
wynika przez dodanie nierówności. Prawa strona
jest mierzalna, a test na półprostych kończy dowód.

W ogólnym przypadku rozdzielamy dziedzinę na mierzalne
zbiory
\[
\begin{aligned}
D_0&=\{|f|<\infty\}\cap\{|g|<\infty\},\\
D_+&=\{f=+\infty,\ g>-\infty\}
      \cup\{g=+\infty,\ f>-\infty\},\\
D_-&=\{f=-\infty,\ g<+\infty\}
      \cup\{g=-\infty,\ f<+\infty\}.
\end{aligned}
\]
W zapisach $|f|<\infty$ i $|g|<\infty$ chodzi o
skończoność wartości; dla punktów z $D_+$ lub $D_-$
suma jest odpowiednio równa $+\infty$ lub $-\infty$.
Założenie wykluczające przeciwne nieskończoności
sprawia, że $D_0,D_+,D_-$ tworzą rozłączny podział
$X$. Na $D_0$ sumę otrzymujemy przez złożenie
mierzalnego odwzorowania $(f,g)|_{D_0}$ z ciągłym
dodawaniem na $\mathbb R^2$. Na dwóch pozostałych
zbiorach funkcja jest stała. Przeciwobrazy
zbiorów borelowskich są zatem mierzalne, bez
odwoływania się do twierdzenia o granicach,
które dopiero zostanie udowodnione.
\end{proof}

\begin{uwaga}[Niezdefiniowane działania]
Dla $f(x)=+\infty$ i $g(x)=-\infty$ formalny zapis
$f(x)+g(x)$ nie określa liczby w $\overline{\mathbb R}$.
Nie można więc pytać o mierzalność takiej \emph{funkcji},
zanim nie ustali się jej wartości na zbiorze punktów
problematycznych. Ustalenie wartości stałej na tym
mierzalnym zbiorze daje jedną z możliwych mierzalnych
wersji rozszerzenia działania.
\end{uwaga}

\section{Granice punktowe funkcji mierzalnych}

\begin{twierdzenie}[Kresy i granice]
\label{thm:pointwise-limits}
Niech $(X,\mathcal F)$ będzie przestrzenią mierzalną,
a $(f_n)_{n\ge1}$ ciągiem mierzalnych funkcji
$f_n:X\to\overline{\mathbb R}$. Wówczas funkcje
\[
\sup_{n\ge1}f_n,\qquad \inf_{n\ge1}f_n,\qquad
\limsup_{n\to\infty}f_n,\qquad
\liminf_{n\to\infty}f_n
\]
są mierzalne. Zbiór
$D=\{x:\lim_{n\to\infty}f_n(x)\text{ istnieje w }
\overline{\mathbb R}\}$ należy do $\mathcal F$;
granica $f(x)=\lim_nf_n(x)$ określona na $D$ jest
mierzalna względem $\mathcal F|_D$.
\end{twierdzenie}
\begin{proof}
Dla dowolnego $a\in\mathbb R$ mamy
$\{\sup_nf_n\le a\}=\bigcap_n\{f_n\le a\}\in\mathcal F$,
a także $\{\inf_nf_n<a\}=\bigcup_n\{f_n<a\}$.
To dowodzi mierzalności supremum i infimum.
Ponieważ
$\limsup_nf_n=\inf_{m\ge1}\sup_{n\ge m}f_n$ i
$\liminf_nf_n=\sup_{m\ge1}\inf_{n\ge m}f_n$,
obie funkcje są mierzalne.
W rozszerzonej prostej granica istnieje wtedy
i tylko wtedy, gdy $\liminf_nf_n=\limsup_nf_n$.
Zbiór równości dwóch mierzalnych funkcji rozszerzonych
jest mierzalny: można na przykład zastosować
ciągłe, ściśle rosnące przekształcenie
$\theta:\overline{\mathbb R}\to[-1,1]$,
$\theta(t)=2\arctan(t)/\pi$ dla $t\in\mathbb R$,
$\theta(\pm\infty)=\pm1$, i następnie badać zbiór,
na którym $|\theta(\limsup_nf_n)-
\theta(\liminf_nf_n)|=0$.
Na $D$ granica jest ograniczeniem mierzalnej funkcji
$\limsup_nf_n$, więc jest mierzalna względem śladu.
\end{proof}

\begin{wniosek}[Monotoniczne ciągi funkcji]\label{idx:06:wniosek:06}
Niech $(X,\mathcal F)$ będzie przestrzenią mierzalną.
Jeżeli $(f_n)_{n\ge1}$ jest niemalejącym ciągiem
mierzalnych funkcji $f_n:X\to[0,\infty]$,
to $f=\sup_nf_n$
jest mierzalna i $f_n(x)\uparrow f(x)$ dla każdego $x$.
Analogicznie, jeżeli $(g_n)_{n\ge1}$ jest malejącym ciągiem
mierzalnych funkcji $g_n:X\to\overline{\mathbb R}$,
to funkcja $\inf_ng_n$ jest mierzalna oraz
$g_n(x)\downarrow\inf_ng_n(x)$ dla każdego $x\in X$.
\end{wniosek}

\begin{przyklad}[Granica punktowa funkcji ciągłych]
Funkcje $f_n(x)=x^n$ na $[0,1]$ są ciągłe,
a więc borelowsko mierzalne. Zbiegają punktowo do
$\chi_{\{1\}}$, która jest mierzalna, choć
nie jest ciągła w punkcie $1$ (dla topologii podprzestrzeni).
Granice punktowe zachowują mierzalność, lecz nie muszą
zachowywać ciągłości.
\end{przyklad}

\section{Funkcje proste}

\begin{definicja}[Funkcja prosta]
\label{def:simple-function}
Funkcję $s:X\to\mathbb R$ nazywamy
\emph{funkcją prostą mierzalną}, jeśli przyjmuje
skończenie wiele wartości i jest mierzalna.
Równoważnie istnieją $m\in\mathbb N$,
liczby $a_1,\ldots,a_m\in\mathbb R$ i zbiory
$A_1,\ldots,A_m\in\mathcal F$ tworzące podział $X$
takie, że
\[
s=\sum_{j=1}^m a_j\chi_{A_j}.
\]
Jeśli wszystkie $a_j\ge0$, funkcję prostą nazywamy
\emph{nieujemną}.
\end{definicja}

\begin{proof}[Dowód równoważności w definicji]
Jeżeli funkcja $s$ przyjmuje skończenie wiele różnych
wartości $a_1,\ldots,a_m$, to jej włókna
$A_j=\{x:s(x)=a_j\}$ są mierzalne i tworzą
podział $X$. Na $A_j$ funkcja ma wartość $a_j$,
co daje podaną reprezentację. Odwrotnie, podana suma
jest mierzalna jako suma skończona mierzalnych funkcji
charakterystycznych i ma zbiór wartości zawarty w
$\{a_1,\ldots,a_m\}$.
\end{proof}

\begin{uwaga}[Po co wymagamy podziału?]
Reprezentacja jako skończona kombinacja funkcji
charakterystycznych nie musi być jednoznaczna, jeśli
zbiory się nakładają. Zawsze można ją zastąpić
reprezentacją kanoniczną, w której $A_j$ są rozłączne
i odpowiadają dokładnie różnym wartościom $s$.
Wówczas wzór na całkę funkcji prostej będzie
niezależny od sposobu zapisania tej funkcji.
\end{uwaga}

\begin{twierdzenie}[Monotoniczna aproksymacja funkcji nieujemnej]
\label{thm:simple-up}
Niech $(X,\mathcal F)$ będzie przestrzenią mierzalną,
a $f:X\to[0,\infty]$ funkcją mierzalną. Dla
$n\ge1$ określmy
\begin{equation}\label{eq:canonical-simple}
s_n(x)=2^{-n}\Big\lfloor 2^n\min\{f(x),n\}\Big\rfloor,
\qquad x\in X.
\end{equation}
Wtedy każda $s_n$ jest nieujemną funkcją prostą mierzalną,
$0\le s_n\le s_{n+1}\le f$ oraz
$s_n(x)\to f(x)$ dla każdego $x\in X$,
przy czym granica może wynosić $+\infty$.
Dla $f(x)<\infty$ i $n>f(x)$ zachodzi
$0\le f(x)-s_n(x)<2^{-n}$.
\end{twierdzenie}
\begin{proof}
Wartość $\min(f,n)$ należy do $[0,n]$.
Po pomnożeniu przez $2^n$ i zastosowaniu części
całkowitej otrzymujemy jedną z liczb całkowitych
$0,1,\ldots,n2^n$. Funkcja $s_n$ przyjmuje zatem
skończenie wiele wartości. Dla $k=0,\ldots,n2^n-1$
zbiór $\{s_n=k2^{-n}\}$ jest opisany nierównościami
$k2^{-n}\le\min(f,n)<(k+1)2^{-n}$;
zbiór odpowiadający ostatniej wartości jest
$\{f\ge n\}$. Wszystkie te zbiory są mierzalne,
więc $s_n$ jest funkcją prostą mierzalną.

Wprowadźmy $u_n(x)=\min(f(x),n)$. Mamy
$u_n\le u_{n+1}$; funkcja
$t\mapsto 2^{-n}\lfloor2^nt\rfloor$ jest niemalejąca,
a ponadto dla każdego $t\ge0$
$2^{-(n+1)}\lfloor2^{n+1}t\rfloor\ge
2^{-n}\lfloor2^nt\rfloor$.
Ostatnia nierówność wynika z tego, że
$\lfloor2z\rfloor\ge2\lfloor z\rfloor$.
Wynika stąd $s_n\le s_{n+1}$. Nierówności
$0\le s_n\le u_n\le f$ wynikają z własności
części całkowitej. Gdy $f(x)$ jest skończona,
przy $n>f(x)$ mamy $u_n(x)=f(x)$ i
$0\le f(x)-s_n(x)<2^{-n}$.
Jeżeli $f(x)=\infty$, to $u_n(x)=n$ i
$s_n(x)=n\to\infty$. To kończy dowód.
\end{proof}

\begin{przyklad}[Aproksymacja z podziałem na poziomy]
Niech $X=[0,1]$, $f(x)=x$ i $n=2$.
Wtedy $s_2(x)=\frac14\lfloor4x\rfloor$ ma wartości
$0,1/4,1/2,3/4,1$ na kolejnych przedziałach
$[0,1/4)$, $[1/4,1/2)$, $[1/2,3/4)$,
$[3/4,1)$ i singletonie $\{1\}$.
Dalsze przybliżenia są coraz dokładniejszymi
schodkami leżącymi poniżej wykresu $f$.
\end{przyklad}

\begin{center}
\begin{tikzpicture}[x=9cm,y=4.0cm]
  \draw[->] (-0.015,0) -- (1.08,0) node[right] {$x$};
  \draw[->] (0,-0.025) -- (0,1.1) node[above] {$y$};
  \draw[blue!70!black,thick] (0,0)--(1,1) node[above left] {$f(x)=x$};
  \draw[red!75!black,thick] (0,0)--(.25,0) (.25,.25)--(.5,.25) (.5,.5)--(.75,.5) (.75,.75)--(1,.75);
  \foreach \x/\y in {.25/0,.5/.25,.75/.5,1/.75} {
    \draw[red!75!black,fill=white] (\x,\y) circle[radius=1.2pt];
  }
  \foreach \x/\y in {.25/.25,.5/.5,.75/.75,1/1} {
    \fill[red!75!black] (\x,\y) circle[radius=1.2pt];
  }
  \draw[dashed,gray!60] (1,0)--(1,1);
  \node[anchor=north] at (0,0) {$0$};
  \node[anchor=north] at (1,0) {$1$};
  \node[red!75!black,anchor=west] at (.42,.13) {$s_2$};
\end{tikzpicture}

\smallskip
\emph{Aproksymacja funkcji $f(x)=x$ na $[0,1]$ przez $s_2$.
Wartości w punktach skoku definiuje wzór \eqref{eq:canonical-simple}.}
\end{center}

\begin{twierdzenie}[Aproksymacja funkcji o wartościach rzeczywistych]
\label{thm:simple-real-approx}
Niech $(X,\mathcal F)$ będzie przestrzenią mierzalną,
a $f:X\to\mathbb R$ mierzalną funkcją skończoną.
Wówczas istnieje ciąg funkcji prostych mierzalnych
$(r_n)$, dla którego $r_n(x)\to f(x)$ dla każdego $x\in X$.
Jeśli $|f(x)|\le M$ dla wszystkich $x\in X$ i pewnego
$M<\infty$, można wybrać $r_n$ tak, by
$\|r_n-f\|_\infty\le2^{-n}$, gdzie
$\|h\|_\infty:=\sup_{x\in X}|h(x)|$.
\end{twierdzenie}
\begin{proof}
Niech $f^+=\max(f,0)$ i $f^-=\max(-f,0)$.
Obie funkcje są mierzalne i nieujemne.
Na mocy twierdzenia \ref{thm:simple-up} istnieją
niemalejące ciągi nieujemnych funkcji prostych
$s_n^+\uparrow f^+$ i $s_n^-\uparrow f^-$.
Funkcje $r_n=s_n^+-s_n^-$ są proste i mierzalne:
przyjmują skończenie wiele wartości i są różnicami
funkcji mierzalnych skończonych. Ponadto
$r_n(x)\to f^+(x)-f^-(x)=f(x)$ dla każdego $x$.
Jeśli $f$ jest ograniczona przez $M$, dla
$n>M$ oba błędy aproksymacji są mniejsze niż
$2^{-n}$, a ponieważ w każdym punkcie co najwyżej
jedna z liczb $f^+(x),f^-(x)$ jest dodatnia,
$|r_n(x)-f(x)|<2^{-n}$. Dla początkowych indeksów
można wybrać dowolne funkcje proste, gdyż nie
wpływa to na istnienie późniejszych przybliżeń;
aby uzyskać podaną nierówność dla wszystkich $n$,
bezpośrednio stosujemy obcięcie i kwantyzację
$r_n(x)=2^{-n}\lfloor2^nf(x)\rfloor$.
Przy ograniczonym $f$ funkcja ta jest prosta i
mierzalna, a dla każdego $x$ zachodzi
$0\le f(x)-r_n(x)<2^{-n}$. Dowodzi to ostatniej tezy.
\end{proof}

\begin{wniosek}[Charakteryzacja mierzalności przez funkcje proste]
\label{cor:meas-simple-characterization}
Niech $(X,\mathcal F)$ będzie przestrzenią mierzalną,
a $f:X\to\mathbb R$ funkcją. Wówczas $f$ jest
mierzalna wtedy i tylko wtedy, gdy istnieje
ciąg mierzalnych funkcji prostych $r_n:X\to\mathbb R$
zbieżny punktowo do $f$.
\end{wniosek}
\begin{proof}
Konieczność daje twierdzenie \ref{thm:simple-real-approx}.
Dostateczność wynika z twierdzenia
\ref{thm:pointwise-limits}, ponieważ punktowa granica
mierzalnych funkcji rzeczywistych jest mierzalna.
\end{proof}

\section{Mierzalność funkcji zespolonych i wektorowych}

\begin{definicja}[Funkcja zespolona mierzalna]\label{idx:06:definicja:11}

Funkcję $f:X\to\mathbb C$ nazywamy mierzalną,
jeśli jest $(\mathcal F,\mathcal B(\mathbb C))$-mierzalna,
gdzie $\mathbb C$ utożsamiamy z $\mathbb R^2$
przez $z\mapsto(\operatorname{Re}z,\operatorname{Im}z)$.
\end{definicja}

\begin{twierdzenie}[Test na częściach rzeczywistej i urojonej]
\label{thm:complex-test}
Niech $(X,\mathcal F)$ będzie przestrzenią mierzalną,
a $f:X\to\mathbb C$. Funkcja $f$ jest mierzalna
wtedy i tylko wtedy, gdy funkcje
$\operatorname{Re}f$ i $\operatorname{Im}f$ są
mierzalne jako funkcje rzeczywiste. W takim przypadku
$\overline f$ i $|f|$ także są mierzalne.
\end{twierdzenie}
\begin{proof}
Pierwsza część jest szczególnym przypadkiem
wniosku \ref{cor:coordinate-measurable} dla $d=2$.
Sprzężenie zespolone i moduł są ciągłe na $\mathbb C$,
a złożenie z funkcją mierzalną pozostaje mierzalne.
\end{proof}

\begin{przyklad}[Funkcje parametryczne]
Jeśli $u,v:X\to\mathbb R$ są mierzalne, to
$f(x)=u(x)+iv(x)$ jest mierzalna.
Nie potrzebujemy do tego żadnej dodatkowej miary
na $X$. Gdy zaś $u,v$ są funkcjami mierzalnymi
względem $\mathcal L_d$, ale niekoniecznie
borelowskimi, również $f$ jest mierzalna względem
$\mathcal L_d$.
\end{przyklad}

\section{Własności prawie wszędzie i uzupełnianie miar}

\begin{definicja}[Równość i zbieżność prawie wszędzie]
\label{def:ae-functions}
Niech $(X,\mathcal F,\mu)$ będzie przestrzenią z miarą.
Funkcje mierzalne $f,g:X\to\mathbb R$ są równe
\emph{$\mu$-prawie wszędzie}, jeśli istnieje
$N\in\mathcal F$ o $\mu(N)=0$ takie, że
$f(x)=g(x)$ dla każdego $x\in X\setminus N$.
Ciąg funkcji mierzalnych $f_n:X\to\mathbb R$
\emph{zbiega $\mu$-prawie wszędzie} do $f:X\to\mathbb R$,
co zapisujemy $f_n\to f$ $\mu$-p.w., jeśli istnieje
taki zbiór $N$, że $f_n(x)\to f(x)$ dla każdego
$x\in X\setminus N$.
\end{definicja}

\begin{uwaga}[Ostrożnie z niezupełną miarą]
Jeśli $\mu$ nie jest zupełna, podzbiór zbioru
mierzalnego miary zero nie musi być mierzalny.
Dlatego definicję wyraziliśmy przez istnienie
\emph{mierzalnego} zbioru zerowego zawierającego
zbiór punktów wyjątkowych. Dla funkcji mierzalnych
sam zbiór $\{f\ne g\}$ jest mierzalny, więc można
równoważnie wymagać $\mu(\{f\ne g\})=0$.
Podobnie dla mierzalnych $f_n,f$ zbiór braku
zbieżności jest mierzalny na mocy twierdzenia
\ref{thm:pointwise-limits}.
\end{uwaga}

\begin{twierdzenie}[Modyfikacje na zbiorach zerowych]
\label{thm:null-modifications}
Niech $(X,\mathcal F,\mu)$ będzie \emph{zupełną}
przestrzenią z miarą. Jeżeli $f:X\to\mathbb R$
jest mierzalna, $g:X\to\mathbb R$ dowolna,
a $\{x:f(x)\ne g(x)\}$ jest zawarty w pewnym
$N\in\mathcal F$ o $\mu(N)=0$, to $g$ jest mierzalna.
Bez założenia zupełności wniosek może być fałszywy.
\end{twierdzenie}
\begin{proof}
W przestrzeni zupełnej każdy podzbiór $N$ należy
do $\mathcal F$. Dla zbioru borelowskiego
$B\subseteq\mathbb R$ zbiory
$g^{-1}(B)\cap(X\setminus N)=f^{-1}(B)\cap(X\setminus N)$
i $g^{-1}(B)\cap N\subseteq N$ są mierzalne.
Ich suma jest równa $g^{-1}(B)$, zatem $g$ jest mierzalna.
Dla kontrprzykładu bierzemy nieuzupełnioną borelowską
miarę Lebesgue'a na $\mathbb R$ oraz podzbiór
$A$ zbioru Cantora, który nie jest borelowski
(wniosek \ref{cor:nonborel-lebesgue}).
Funkcja $f\equiv0$ jest mierzalna, ale
$g=\chi_A$ nie jest borelowsko mierzalna.
Jednocześnie $f=g$ poza borelowskim zbiorem
Cantora, którego miara wynosi zero.
\end{proof}

\begin{wniosek}[Wersje mierzalne funkcji określonych p.w.]\label{idx:06:wniosek:15}

Niech $(X,\mathcal F,\mu)$ będzie przestrzenią z miarą,
a $f_n:X\to\mathbb R$ mierzalnymi funkcjami.
Jeżeli ciąg $(f_n(x))$ jest zbieżny dla
$\mu$-prawie wszystkich $x$, to istnieje
mierzalna funkcja $f:X\to\mathbb R$, taka że
$f_n\to f$ $\mu$-p.w.
\end{wniosek}
\begin{proof}
Z twierdzenia \ref{thm:pointwise-limits} zbiór $D$
punktów, w których granica jest rzeczywista i
skończona, jest mierzalny, a granica na nim jest mierzalna.
Z założenia $X\setminus D$ ma miarę zero.
Definiujemy $f(x)=\lim_nf_n(x)$ na $D$ oraz
$f(x)=0$ na $X\setminus D$. Funkcja jest mierzalna,
gdyż dla borelowskiego $B$ jej przeciwobraz jest
sumą mierzalnego przeciwobrazu na $D$ i,
ewentualnie, całego $X\setminus D$.
\end{proof}

\begin{uwaga}[Zbieżność prawie wszędzie nie zapewnia mierzalności dowolnej wersji granicy]
\label{rem:noncomplete-ae-limit-6}
Jeśli nie wymagamy mierzalności funkcji granicznej,
z założenia $f_n\to f$ poza zbiorem zerowym nie można
wnioskować, że $f$ jest mierzalna. Na przestrzeni
$(\mathbb R,\mathcal B(\mathbb R),\lambda_1|_{\mathcal B})$
wybierzmy nieborelowski podzbiór $A$ zbioru Cantora.
Funkcje $f_n\equiv0$ są borelowskie i zbiegają do
$f=\chi_A$ poza borelowskim zbiorem Cantora miary zero,
lecz $f$ nie jest mierzalna względem $\mathcal B(\mathbb R)$.
Wniosek o mierzalnej wersji granicy gwarantuje istnienie
\emph{pewnego} mierzalnego przedstawiciela; nie oznacza,
że każdy wybór wartości na zbiorze wyjątkowym jest
mierzalny. W przestrzeni zupełnej takie modyfikacje są
mierzalne na mocy twierdzenia~\ref{thm:null-modifications}.
\end{uwaga}

\section{Zbieżność według miary: narzędzia potrzebne do całkowania}
Minimalna definicja i twierdzenie o podciągu pojawiają się tutaj,
ponieważ wykorzystamy je już w rozdziale 8; systematyczne porównanie
wszystkich rodzajów zbieżności pozostawiamy rozdziałowi 9.


\begin{definicja}[Zbieżność według miary]
\label{def:convergence-measure}
Niech $(X,\mathcal F,\mu)$ będzie przestrzenią z miarą,
a $f_n,f:X\to\mathbb R$ mierzalnymi funkcjami.
Mówimy, że $f_n$ \emph{zbiega do $f$ według miary $\mu$},
i piszemy $f_n\xrightarrow{\mu}f$, jeśli
\[
\forall\varepsilon>0\qquad
\mu\bigl(\{x\in X:|f_n(x)-f(x)|>\varepsilon\}\bigr)
\xrightarrow[n\to\infty]{}0.
\]
Definicja nie wymaga $\mu(X)<\infty$ ani
$\sigma$-skończoności.
\end{definicja}

\begin{twierdzenie}[Prawie wszędzie implikuje zbieżność według miary przy mierze skończonej]
\label{thm:ae-to-measure-finite}
Jeśli $(X,\mathcal F,\mu)$ jest przestrzenią
z $\mu(X)<\infty$, a $f_n,f:X\to\mathbb R$ są
mierzalne i $f_n\to f$ $\mu$-prawie wszędzie,
to $f_n\xrightarrow{\mu}f$.
\end{twierdzenie}
\begin{proof}
Ustalmy $\varepsilon>0$ i niech
$E_n=\bigcup_{k\ge n}\{|f_k-f|>\varepsilon\}$.
Wtedy $E_n\in\mathcal F$, $E_{n+1}\subseteq E_n$,
a $\bigcap_nE_n$ składa się z punktów, gdzie
nierówność $|f_k-f|>\varepsilon$ zachodzi
nieskończenie wiele razy. Zbieżność prawie wszędzie
gwarantuje $\mu(\bigcap_nE_n)=0$.
Ponieważ $\mu(E_1)\le\mu(X)<\infty$,
ciągłość miary z góry (twierdzenie
\ref{thm:measure-continuity}) daje $\mu(E_n)\to0$.
Wreszcie $\{|f_n-f|>\varepsilon\}\subseteq E_n$,
więc z monotoniczności $\mu(\{|f_n-f|>\varepsilon\})\to0$.
\end{proof}

\begin{przyklad}[Dlaczego skończoność miary jest ważna]
\label{ex:aenotmeasure}
Na $(\mathbb R,\mathcal L_1,\lambda_1)$ niech
$f_n=\chi_{[n,n+1]}$. Dla każdego ustalonego
$x\in\mathbb R$ mamy $f_n(x)=0$ dla wszystkich
dostatecznie dużych $n$, a więc $f_n\to0$ wszędzie.
Jednak dla $\varepsilon=1/2$ zachodzi
$\lambda_1(\{|f_n|>1/2\})=1$ dla wszystkich $n$.
Nie ma więc zbieżności według miary na całej $\mathbb R$.
\end{przyklad}

\begin{twierdzenie}[Podciąg zbieżny prawie wszędzie]
\label{thm:measure-subsequence}
Niech $(X,\mathcal F,\mu)$ będzie dowolną przestrzenią
z miarą, a $f_n,f:X\to\mathbb R$ mierzalnymi funkcjami.
Jeśli $f_n\xrightarrow{\mu}f$, to istnieje ściśle
rosnący ciąg indeksów $(n_k)_{k\ge1}$ taki, że
$f_{n_k}\to f$ $\mu$-prawie wszędzie.
\end{twierdzenie}
\begin{proof}
Z definicji zbieżności według miary wybieramy
indukcyjnie $n_k>n_{k-1}$ tak duże, aby
\[
\mu(A_k)<2^{-k},\qquad
A_k=\{|f_{n_k}-f|>2^{-k}\}.
\]
Wtedy $\sum_{k\ge1}\mu(A_k)\le1<\infty$.
Pierwsza część lematu Borela--Cantellego
(twierdzenie \ref{thm:borel-cantelli-1})
daje $\mu(\limsup_k A_k)=0$.
Poza tym zbiorem punkt $x$ należy do zaledwie
skończenie wielu $A_k$; dla dostatecznie dużych $k$
mamy więc $|f_{n_k}(x)-f(x)|\le2^{-k}$.
Prawa strona zbiega do zera, a zatem
$f_{n_k}(x)\to f(x)$.
\end{proof}

\begin{uwaga}[Dlaczego nie potrzebujemy skończoności miary?]
Wybór podciągu opiera się wyłącznie na sumowalności
$\sum_k\mu(A_k)$, a lemat Borela--Cantellego nie
zakłada $\mu(X)<\infty$. Jest to zasadniczo inny
argument niż ciągłość miary z góry zastosowana w
twierdzeniu~\ref{thm:ae-to-measure-finite}.
\end{uwaga}

\begin{wniosek}[Jednoznaczność granicy według miary]
\label{cor:measure-limit-unique}
Niech $(X,\mathcal F,\mu)$ będzie przestrzenią z miarą,
a $(f_n)_{n\ge1}$ ciągiem mierzalnych funkcji
$f_n:X\to\mathbb R$. Jeżeli mierzalne funkcje
$f,g:X\to\mathbb R$ spełniają $f_n\xrightarrow{\mu}f$
oraz $f_n\xrightarrow{\mu}g$, to $f=g$ $\mu$-prawie wszędzie.
\end{wniosek}
\begin{proof}
Dla $\varepsilon>0$ z nierówności trójkąta
$\{|f-g|>\varepsilon\}\subseteq
\{|f-f_n|>\varepsilon/2\}\cup
\{|f_n-g|>\varepsilon/2\}$.
Z podaddytywności i przejścia do granicy
otrzymujemy $\mu(\{|f-g|>\varepsilon\})=0$.
Ponieważ $\{f\ne g\}=\bigcup_{m\ge1}\{|f-g|>1/m\}$,
zbiór ten ma miarę zero.
\end{proof}

\begin{przyklad}[Zbieżność według miary bez zbieżności p.w.]
\label{ex:typewriter}
Na $X=[0,1)$ z miarą Lebesgue'a, dla każdego $k\ge0$
i $j=0,\ldots,2^k-1$, kładziemy
$f_{2^k+j}=\chi_{[j2^{-k},(j+1)2^{-k})}$.
Każdy indeks $n\ge1$ występuje w tym wzorze dokładnie
raz (blok indeksów $2^k,\ldots,2^{k+1}-1$).
Dla $n=2^k+j$ mamy
$\lambda_1(\{|f_n|>1/2\})=2^{-k}\to0$,
więc $f_n\to0$ według miary.
Z drugiej strony dla każdego $x\in[0,1)$ i każdego
$k$ dokładnie jeden z przedziałów poziomu $k$
zawiera $x$. Stąd w każdym bloku istnieje jeden indeks
$n$ z $f_n(x)=1$ oraz (dla $k\ge1$) inny indeks
z $f_n(x)=0$. Ciąg $f_n(x)$ nie ma granicy w żadnym
punkcie $[0,1)$.
\end{przyklad}

\begin{twierdzenie}[Lokalna zbieżność według miary]
\label{thm:local-convergence-measure}
Niech $(X,\mathcal F,\mu)$ będzie przestrzenią z miarą,
a $f_n,f:X\to\mathbb R$ funkcjami mierzalnymi,
przy czym $f_n\to f$ $\mu$-prawie wszędzie.
Wówczas dla każdego $A\in\mathcal F$ z $\mu(A)<\infty$
i każdego $\varepsilon>0$ zachodzi
\[
\mu\bigl(A\cap\{|f_n-f|>\varepsilon\}\bigr)\longrightarrow0.
\]
\end{twierdzenie}
\begin{proof}
Rozpatrujemy przestrzeń
$(A,\mathcal F|_A,\mu|_A)$ o skończonej mierze,
a następnie stosujemy twierdzenie
\ref{thm:ae-to-measure-finite} do funkcji ograniczonych
do $A$.
\end{proof}

\begin{uwaga}[Co jest, a co nie jest implikacją?]
Zbieżność prawie wszędzie implikuje zbieżność
według miary przy $\mu(X)<\infty$, lecz nie musi
implikować jej dla $\mu(X)=\infty$.
Zbieżność według miary zawsze pozwala wybrać podciąg
zbieżny prawie wszędzie, ale nie wymusza zbieżności
prawie wszędzie całego ciągu. Kontrprzykłady powyżej
pokazują niezależność obu problemów. Twierdzenia
Jegorowa i Riesza o podciągach oraz pełne zestawienie
rodzajów zbieżności rozwinie rozdział 9.
\end{uwaga}

\section{Przygotowanie do definicji całki Lebesgue'a}

Skończone kombinacje funkcji charakterystycznych są
pomostem pomiędzy mierzeniem zbiorów a całkowaniem.
Dla nieujemnej funkcji prostej $s=\sum_{j=1}^m
a_j\chi_{A_j}$, gdzie $A_j$ tworzą mierzalny
podział $X$, naturalny kandydat na całkę to
\[
\sum_{j=1}^m a_j\mu(A_j).
\]
Jeśli $a_j=0$ i $\mu(A_j)=\infty$, przyjmujemy
w tym nieujemnym wyrażeniu konwencję $0\cdot\infty=0$.
W następnym rozdziale dowiedziemy, że liczba ta
nie zależy od reprezentacji $s$, po czym określimy
całkę dowolnej funkcji nieujemnej jako supremum
po wszystkich prostych funkcjach mierzalnych
mniejszych od niej.

W ten sposób twierdzenie \ref{thm:simple-up}
ma fundamentalną rolę: zapewnia, że zbiór takich
przybliżeń jest dostatecznie bogaty, aby odtworzyć
wartości dowolnej funkcji mierzalnej.
